%%%%%%%%%%%%%%%%%%%%%%%%%%%%%%%%%%%%%%%%%%%%%%%%%%%%%%%%%
%%%%%%%%%%%%%%%%%         FUTUROS DESA              %%%%%%%%%%%%%%%%%%%%%%
%%%%%%%%%%%%%%%%%%%%%%%%%%%%%%%%%%%%%%%%%%%%%%%%%%%%%%%%%


\chapter{Futuros desarrollos}
\label{chap:future}

A raíz de la realización de la presentes Tesis Doctoral, se presentan varias líneas de investigación que serían interesantes en futuros desarrollos para explotar la información que nos brinda el microbioma.


\textbf{Integración de plataformas de secuenciación para facilitar comparaciones y combinaciones de datos para obtener resultados más robustos y coherentes}


Como se ha expuesto en la presente Tesis Doctoral, cada tipo de secuenciación del microbioma presenta sus ventajas. Mientras que el 16S rRNA es más asequible, rápido y presenta menos contaminación del ADN del anfitrión, el WGS presenta mayor resolución taxonómica y no presenta un sesgo debido a la selección de primers.

En los últimos años ha surgido la idea de si los datos provenientes de ambos tipos de secuenciación se podrían integrar eficazmente, ya que esto habilitaría diversas posibilidades, como incrementar el número de muestras accesibles de estudios del microbioma y comparar los resultados obtenidos en uno de los tipos de secuenciación con los obtenidos en la plataforma complementaria. Esto se traduciría en un mayor poder estadístico y en una mayor reproducibilidad en el descubrimiento de posibles biomarcadores en el campo del microbioma.

Ante esta necesidad de armonizar datos provenientes de diferentes tecnologías de secuenciación, recientemente se ha desarrollado la base de datos GreenGenes2, que, a futuro, puede constituir una herramienta importante para la integración de secuenciación 16S rRNA y WGS. GreenGenes2 \cite{mcdonald2024greengenes2} se basa en la unificación de secuencias derivadas tanto de la secuenciación de 16S rRNA como de WGS, construyendo un árbol filogenético de referencia que integra información de ambas plataformas en un solo marco taxonómico coherente. Mediante su estrategia de homogeneización, esta base de datos permite que los taxones obtenidos en estudios basados en 16S rRNA se puedan comparar de manera directa con aquellos obtenidos mediante WGS, facilitando así la integración de conjuntos de datos heterogéneos y la realización de meta-análisis a gran escala \cite{mcdonald2024greengenes2}.

Sin embargo, la variabilidad técnica proveniente de los diferentes protocolos de secuenciación o de las diferencias de resolución taxonómica es un gran obstáculo para una integración efectiva. En esta línea, se ha realizado un primer trabajo en el que se ha explorado el comportamiento de varios métodos de reducción de \textit{batch-effect} para integrar ambas plataformas \cite{fernandez2024cross}. En este trabajo se observa cómo estos métodos pueden llegar a corregir de manera eficiente las diferencias entre ambas plataformas. Sin embargo, no es concluyente si la calidad o la integridad de la señal biológica se mantiene después de aplicar dichos métodos.

En resumen, el desarrollo de metodologías y tecnologías que permitan la integración de datos de secuenciación 16S y WGS, así como la adopción de herramientas unificadoras como GreenGenes2, serán claves para avanzar en el análisis computacional del microbioma humano y en la generación de biomarcadores clínicos. Estos esfuerzos fortalecerán la capacidad de realizar estudios multicohorte integrados, optimizarán la comparación entre diferentes plataformas de secuenciación y contribuirán a la construcción de un marco analítico robusto que sea aplicable en entornos clínicos.


\textbf{Integración con otros datos ómicos procedentes del microbioma para obtener una visión más detallada de la actividad y funcionalidad del microbioma}


El análisis composicional del microbioma ha permitido explicar un porcentaje significativo de la patogénesis de ciertas enfermedades. Sin embargo, como se ha comentado en esta Tesis Doctoral, el enfoque metagenómico por sí solo no refleja toda la actividad real de estos microorganismos. Es aquí donde la integración con datos procedentes de metatranscriptómica (que se enfoca en qué genes se están expresando), metaproteómica (que cuantifica proteínas expresadas) y metabolómica (que analiza el perfil de metabolitos) resulta fundamental para revelar el estado funcional real del microbioma.

Los datos multiómicos que ya se están generando permiten medir la actividad real del microbioma, evaluando rutas como la producción de ácidos grasos de cadena corta, el metabolismo de ácidos biliares o la biosíntesis de vitaminas; definir grupos funcionales y relaciones entre microbios con impacto clínico; vincular genes y cepas con fenotipos y fallos de tratamiento; relacionar metabolitos y proteínas microbianas con medidas del huésped (marcadores inmunes y clínicos) para priorizar mecanismos; construir modelos predictivos interpretables para estratificar pacientes y anticipar la respuesta terapéutica, y analizar la evolución temporal para detectar cambios de estado y señales tempranas de descompensación. En conjunto, pasar de las abundancias taxonómicas a las funciones activas y su efecto en el huésped mejora la explicación de la enfermedad y acerca estos análisis a la aplicación clínica.

En los próximos años, la integración de datos metagenómicos con otras capas ómicas permitirá pasar de descripciones composicionales a perfiles funcionales y dinámicos del microbioma en contextos fisiopatológicos. Se adoptarán marcos de integración y se desarrollarán flujos de trabajo bioinformáticos reproducibles y robustos, con el objetivo de establecer biomarcadores clínicos para la estratificación de pacientes y orientar intervenciones terapéuticas basadas en la modulación del microbioma. Este avance hacia enfoques multiómicos integrados se perfila como un paso decisivo que, a medio y largo plazo, mejorará el diagnóstico y optimizará el tratamiento de enfermedades complejas.
    
    
\textbf{Integración microbioma-anfitrión para entender mejor las interacciones biológicas y sus implicaciones en procesos patológicos}


Llegados a este punto, todavía es comprensible pensar que abordar únicamente los datos ómicos derivados del microbioma (metagenoma, metatranscriptoma, metaproteoma y metaboloma) no resulta suficiente para descifrar en su totalidad los mecanismos biológicos subyacentes a determinadas patologías. En este contexto, la integración de datos procedentes tanto del microbioma como del anfitrión (genómica, transcriptómica, proteómica, entre otros) se presenta como una estrategia fundamental para obtener una visión sistémica de los procesos patológicos. Esta aproximación, que combina los múltiples niveles de información de ambas partes, permite detectar interacciones críticas que, de otro modo, podrían pasar desapercibidas.

La justificación de integrar datos del anfitrión radica en que el fenotipo de una enfermedad se encuentra en la intersección entre la información hereditaria del paciente y las influencias ambientales, de las cuales el microbioma es uno de los componentes más dinámicos y moduladores. Diversos estudios han demostrado que el genoma del anfitrión y su expresión fenotípica afectan directamente la composición y el funcionamiento del microbioma; por ejemplo, determinadas variantes genéticas están estrechamente relacionadas con la abundancia de taxones microbianos específicos, lo que repercute en procesos como la respuesta inmune, el metabolismo y la inflamación. Asimismo, las modificaciones inducidas en el transcriptoma y proteoma del huésped, ya sea por respuestas inmunitarias o por cambios metabólicos, pueden modular la actividad microbiana y, a su vez, influir en el desarrollo de patologías complejas \cite{duan2025advances}. Para abordar estas interacciones bidireccionales, es necesario el desarrollo de enfoques computacionales que permitan la integración simultánea de diferentes plataformas.

Un ejemplo revelador corresponde a los estudios en pacientes con cáncer, en los que se ha observado que la composición del microbioma puede influir en la eficacia de la inmunoterapia. La integración de perfiles metagenómicos con análisis del transcriptoma y proteoma tumoral ha permitido identificar biomarcadores que predicen la respuesta al tratamiento, lo que abre la puerta a intervenciones que combinan escenarios de microbiota dirigida con terapias convencionales \cite{goyal2025understanding}. Otro campo en el que la integración de datos del anfitrión y el microbioma ha mostrado resultados prometedores es el de la sepsis. Estudios recientes han empleado enfoques multiómicos para combinar datos de la expresión génica del huésped, perfiles proteómicos y datos metagenómicos, lo cual ha permitido distinguir entre infecciones bacterianas y virales en tiempo real, identificando biomarcadores que pueden predecir de manera temprana la progresión hacia cuadros de sepsis grave \cite{lu2025multi}.

En conjunto, la integración microbioma-huésped constituirá un eje clave de los futuros desarrollos: mejorará la identificación y validación de biomarcadores clínicos, permitirá modelos mecanísticos más precisos y favorecerá estrategias terapéuticas personalizadas que modulen tanto la respuesta del huésped como la composición y la función del microbioma. Esta línea de trabajo acercará los análisis multiómicos a la práctica clínica y optimizará la estratificación y el tratamiento de enfermedades complejas.


\textbf{Análisis del papel del microbioma en la resistencia a tratamientos terapéuticos} 


La compleja interacción entre el microbioma humano y los tratamientos farmacológicos constituye un área de investigación en rápida expansión, en la que se ha evidenciado que tanto el microbioma intestinal como el microbioma asociado a tejidos (por ejemplo, el microambiente tumoral) pueden actuar como barreras que impiden que algunos fármacos alcancen sus objetivos terapéuticos. El concepto de farmacomicrobiómica abarca la capacidad del microbioma para modificar, inactivar o incluso reactivar los fármacos mediante una serie de mecanismos que incluyen, entre otros, la captación de los compuestos, la secreción de metabolitos que alteran su estructura química y la modulación de la respuesta inmunitaria del huésped.

El microbioma puede actuar como barrera frente a fármacos a través de diferentes mecanismos directos e indirectos. Por ejemplo, en estudios experimentales se ha evidenciado que ciertas bacterias poseen la capacidad de secuestrar fármacos, resultando en menores concentraciones del compuesto en el lumen intestinal. Esta captación puede ser especialmente crítica en el caso de medicamentos con ventana terapéutica estrecha, ya que pequeñas variaciones en la biodisponibilidad pueden marcar la diferencia entre eficacia y toxicidad \cite{scher2020pharmacomicrobiomics}. Otro mecanismo importante es la secreción de metabolitos que pueden modificar la estructura química de los fármacos. Se han documentado casos en los cuales la interacción entre metabolitos bacterianos y fármacos cardiovasculares genera una notable variabilidad en la respuesta terapéutica \cite{steiner2022role}. Además, el microbioma influye indirectamente en la resistencia a través de la modulación de las rutas enzimáticas del huésped. Por ejemplo, los metabolitos microbianos pueden competir con los fármacos por sitios activos enzimáticos. Este tipo de interacciones evidencia que, más allá de la captación y secreción directa, el microbioma juega un rol subyacente en la regulación de la respuesta del sistema de metabolización del huésped, actuando en sinergia con la genética del paciente para determinar el destino del fármaco \cite{verdegaal2024integrating}.

La relevancia clínica del estudio de la resistencia microbiana a tratamientos es indudable. La capacidad del microbioma para actuar como barrera frente a tratamientos farmacológicos tiene implicaciones clínicas de alta relevancia. Por un lado, la variabilidad interindividual en la composición y función del microbioma puede explicar diferencias significativas en la eficacia y toxicidad de fármacos entre pacientes, lo que plantea un reto en el diseño de terapias universales. La identificación de biomarcadores microbianos predictivos de respuesta o resistencia se torna, por tanto, esencial para personalizar los tratamientos y evitar fracasos terapéuticos \cite{dempsey2019microbiome}. Por otro lado, la presencia de un microbioma tumoral disfuncional podría obstaculizar la acción de terapias dirigidas, como los inhibidores de puntos de control inmunitarios, al generar un microambiente que favorezca la inactivación de fármacos o la supresión de la respuesta inmune. En este sentido, intervenciones que modulen la composición microbiana representan una vía prometedora para aumentar la eficacia terapéutica y reducir la toxicidad asociada \cite{garajova2021role}.

En conclusión, el microbioma se posiciona como un factor importante en la modulación de la respuesta a tratamientos farmacológicos, actuando como barrera que impide que los fármacos alcancen su objetivo terapéutico. Estos mecanismos contribuyen a la variabilidad en la eficacia y toxicidad de los tratamientos, lo que plantea desafíos significativos en el manejo clínico y la personalización de terapias. La integración de datos ómicos del microbioma y del anfitrión, a través de herramientas computacionales y enfoques multiómicos, ofrece una oportunidad única para identificar biomarcadores predictivos y diseñar intervenciones que modulen la barrera microbiana. A futuro, la combinación de estrategias de modelado, ensayos clínicos longitudinales y terapias moduladoras del microbioma promete transformar el enfoque terapéutico, permitiendo optimizar la eficacia de los tratamientos y reducir la incidencia de resistencias adquiridas.